%% ---------------------------------------------------------------------
%% 03 - KẾT QUẢ
%% Bảng: đặt \caption phía TRÊN bảng, chữ 9 pt, nhãn "Bảng 1:" in đậm.
%% Hình: đặt \caption phía DƯỚI hình, caption thụt vào ~9 pt như mẫu.
%% Hình minh hoạ dưới đây vẽ bằng TikZ để không cần file ảnh bên ngoài;
%% khi có hình thật, thay bằng \includegraphics[width=\linewidth]{...}
%% ---------------------------------------------------------------------

\section{Kết quả}

\subsection{Ma trận luật của bộ giám sát}

Bảng~\ref{tab:rules} liệt kê chín luật mờ dùng trong bộ giám sát. Mỗi ô là
mức can thiệp tương ứng với cặp (khoảng cách, tốc độ tương đối).

%% Bảng đặt \caption phía trên, chữ 9 pt; nội dung bảng 7 pt như mẫu.
%% Bảng canh trái theo lề cột (KHÔNG \centering) và được phép tràn nhẹ
%% sang phải nếu rộng hơn cột, đúng như bản in của tạp chí.
\begin{table}[!htbp]
  \caption{Ma trận hành động của bộ giám sát logic mờ.
           [Nguồn: Nhóm tác giả]}
  \label{tab:rules}
  \raggedright
  {\fontsize{7}{9.3}\selectfont
  \begin{tabular}{@{}lccc@{}}
    \toprule
    & \multicolumn{3}{c}{\textbf{Độ tin cậy bản đồ $V$}} \\
    \textbf{Pin $B$} & Tin cậy & Không chắc chắn & Không hợp lệ \\
    \midrule
    Cao        & C   & R   & S+R \\
    Trung bình & C   & R   & S+R \\
    Thấp       & S+R & S+R & S+R \\
    \bottomrule
  \end{tabular}}
\end{table}
% Fix potential mis-order of Bảng references: ensure any further mention of Bảng 1
% refers to the already-introduced table above.
Kết quả được tổng hợp ở Bảng~\ref{tab:rules} ở trên và có thể tham khảo lại tại Bảng~\ref{tab:rules}.

\subsection{Kiến trúc ba lớp}

Hình~\ref{fig:pipeline} minh hoạ luồng xử lý: dữ liệu LiDAR được gom cụm
thành vật cản, xấp xỉ bằng $\beta$-complexes, sinh mạng Voronoi có trọng
số, rồi chuyển sang bộ lấy mẫu có hướng dẫn trước khi bộ giám sát mờ quyết
định có lập lại kế hoạch hay không.

\begin{figure}[!htbp]
  \centering
  \begin{tikzpicture}[
      node distance=7mm,
      layer/.style={draw=black!70, rounded corners=2pt, fill=jstblue,
                    text width=138pt, align=center, font=\footnotesize,
                    inner sep=4pt, minimum height=15pt},
      arr/.style={-latex, black!70, line width=0.5pt}]
    \node[layer] (lidar) {LiDAR 2D $\rightarrow$ gom cụm vật cản};
    \node[layer, below=of lidar] (geom) {Voronoi có trọng số
        $+$ $\beta$-complexes};
    \node[layer, below=of geom] (plan) {RRT* có hướng dẫn};
    \node[layer, below=of plan] (fuzzy) {Giám sát logic mờ};
    \draw[arr] (lidar) -- (geom);
    \draw[arr] (geom) -- (plan);
    \draw[arr] (plan) -- (fuzzy);
    \draw[arr, dashed] (fuzzy.east) -- ++(10pt,0) |- (plan.east);
  \end{tikzpicture}
  \caption{Kiến trúc ba lớp của khung lập kế hoạch lai ghép.
           Đường nét đứt là vòng phản hồi khi bộ giám sát yêu cầu lập lại
           kế hoạch. [Nguồn: Nhóm tác giả]}
  \label{fig:pipeline}
\end{figure}

\subsection{Chỉ số đánh giá}

Ba nhóm chỉ số được ghi nhận: thời gian lập kế hoạch, độ dài quỹ đạo và số
lần lập lại kế hoạch trên mỗi mét di chuyển. Mỗi cấu hình được chạy 30 lần
với cùng bản đồ và cùng phân bố vật cản để bảo đảm ý nghĩa thống kê.
